\section{FLORES-200：先建评估基础设施}

有了真实需求，下一步不是立刻训练，而是先回答“怎样知道系统真的进步”。低资源研究常见困境是训练集、验证集和指标都缺失；没有共同、可比、人工翻译的评测集，模型改进无法跨语言验证。FLORES 系列就是为这道基础设施缺口而建。

\subsection{Evaluation-first 的研究顺序}

本节解释为什么课堂先讲 evaluation、后讲 training。课堂明确把 evaluation datasets 放在 training datasets 之前。这个顺序表达了一条工程纪律：若成功条件、测试分布和失败类型尚未定义，训练迭代只会优化最方便测量的代理目标。

\lecturefigure{slide-09-evaluation-datasets.jpg}{NLLB 先讨论评测数据，再讨论训练数据}{Stanford 课堂视频 00:07:24}

\begin{knowledgebox}{读图：分节页也在表达优先级}
这张 slide 几乎没有技术细节，却给出课程组织逻辑：benchmark 不是模型完成后的附属物，而是训练方向的坐标系。FLORES-200 先定义跨语言可比性，后续 data mining、MoE 和 curriculum 才能在同一尺度上判断得失。
\end{knowledgebox}

\teachervoice{Fan 说自己要“先谈 evaluation data，因为把 evaluation 做准极其重要”。这句话应当读成研究方法：不要在没有可靠测试集时，用训练 loss 或少数高资源语言的结果代表全局。}

\subsection{从 FLORES 到 FLORES-200}

接下来需要理解 benchmark 如何逐代扩展。FLORES 最早围绕少量低资源语言建立评测，随后扩展到 FLORES-101，NLLB 再把覆盖推进到 FLORES-200。课堂提到名字源自 Facebook Low Resource，即使公司改名 Meta，项目仍保留 FLORES 这个名称。

\lecturefigure{slide-10-flores-101.jpg}{FLORES-101 为大规模低资源评测建立前身}{Stanford 课堂视频 00:08:12}

\begin{knowledgebox}{读图：benchmark 的价值来自共同源句与人工翻译}
FLORES 让多种语言共享一批语义对应的句子，因此可以比较 into-English、out-of-English 和 non-English translation。它不是抓取现成网页再随机切分，而是用专业翻译和质量流程生产更稳定的评测锚点。
\end{knowledgebox}

\begin{importantbox}{翻译方向为什么会组合爆炸}
若有 $N$ 种语言，并把 $A\rightarrow B$ 与 $B\rightarrow A$ 视为不同方向，则最多存在
\[
D=N(N-1)
\]
个有向翻译方向。$N$ 表示语言数量，$D$ 表示不含同语种复制的方向数。对 $N=200$，$D=39{,}800$；这解释了课堂所说的四万余方向，也解释了为什么训练数据不可能对每个方向都同样充分。
\end{importantbox}

\subsection{FLORES 的四个设计属性}

本节进一步把“有一个评测集”拆成具体设计。课堂用 progressive bullets 展示 FLORES 的设计选择，最终状态包含四点：聚焦低资源语言、支持 many-to-many、覆盖多样主题与领域、保留 document-level context。这些属性分别应对覆盖偏差、方向偏差、领域偏差和孤立句偏差。

\lecturefigure{slide-11-flores-properties.jpg}{FLORES-200 的低资源、多对多、跨领域与文档级设计}{Stanford 课堂视频 00:09:57}

\begin{knowledgebox}{读图：四个属性对应四种 benchmark 失真}
只选高资源语言会高估系统成熟度；只测 English-centric directions 会隐藏 non-English 退化；只测单一领域会把领域匹配误写成通用能力；只测孤立句会忽略代词、连贯性和上下文歧义。FLORES 不能消灭所有偏差，但它让这些失真显著减小并可被讨论。
\end{knowledgebox}

\begin{warningbox}{Many-to-many evaluation 不等于 many-to-many supervision}
评测集可以构造所有语言对的对照，不代表训练阶段拥有所有方向的平行语料。课堂 Q\&A 明确说，NLLB 没有挖掘全部 $200\times200$ 方向；许多方向依赖共享多语表示产生 zero-shot transfer。
\end{warningbox}

\subsection{专业翻译、自动检查与人工复审}

FLORES 的生产流程从语言标准协商开始，然后把文档交给译者，经过自动检查、人工 reviewer 与 post-editing。若质量低于门槛，样本回到修订环节。这个循环说明高质量评测数据不是一次采购，而是需要语言标准、工具检查与专业判断共同约束。

\lecturefigure{slide-12-flores-collection-pipeline.jpg}{FLORES 从语言标准协商到复审与后编辑的流程}{Stanford 课堂视频 00:11:30}

\begin{knowledgebox}{读图：箭头中的回路比最后的完成节点更重要}
先确定 script、拼写和变体标准，译者才知道目标文本应采用哪种形式；自动检查可以发现漏译、格式和异常字符，却不能独立判断语义自然度；reviewer 与 post-editing 负责把失败样本送回。图中的 `Quality > 90%` 是该流程的课堂展示条件，不应外推为所有语言项目的固定治理阈值。
\end{knowledgebox}

\teachervoice{老师强调，真正困难的第一步往往是 alignment on language standards。若社区同时使用多种脚本、拼写和本地变体，团队必须先决定评测对象是什么，不能等模型训练完再补标签。}

\subsection{翻译者、标准、脚本与地方变体}

最后回到评测集的生产约束。最后一张评测数据 slide 把收集难点压缩为三类：找到合适译者，协调语言标准，以及处理多脚本与地方变体。它们看似是“数据运营”，实际上直接决定 benchmark 是否测到了目标语言，还是只测到一种机构化书写形式。

\lecturefigure{slide-13-flores-collection-challenges.jpg}{FLORES 数据收集中的译者与语言标准挑战}{Stanford 课堂视频 00:12:54}

\begin{knowledgebox}{读图：每个 bullet 都会传播到模型与指标}
译者招募不足会扩大方差；标准选择不当会把合法变体判成错误；脚本混杂会破坏 tokenization 与 language identification；地方变体缺失会让模型在 benchmark 上很好，却无法服务真实社群。数据标准不是中立容器，而是整个系统的 measurement model。
\end{knowledgebox}

\begin{warningbox}{Translationese 与 English-centric source}
人工翻译评测集仍可能带有 translationese，即目标句保留源语言结构；若源内容主要来自英语语境，也可能缺少本地文化主题。高质量人工翻译显著改善可比性，但并不自动等于原生语料分布。
\end{warningbox}

\subsection{本章小结}

FLORES-200 先为两百语言建立共同评估坐标，再让模型研究可比较。其价值来自低资源覆盖、多对多、跨领域、文档上下文和专业翻译流程；局限则包括标准选择、translationese 与源内容偏差。

